\chapter{Conjunto de datos}
\label{cap:datos}

\section{Origen y estructura}

Los datos provienen de transacciones reales de comercio electrónico provistas por Vesta
Corporation. Se componen de dos tablas unidas por \texttt{TransactionID}
(\cref{tab:datos}): una tabla de \emph{transacción} (590\,540 filas $\times$ 394 columnas en
entrenamiento) y una tabla de \emph{identidad} (144\,233 filas $\times$ 40 columnas), disponible
solo para el $\sim$24\,\% de las transacciones.

\begin{table}[ht]
  \centering
  \caption{Resumen de los archivos del conjunto de datos.}
  \label{tab:datos}
  \input{\tabdir tabla_datos}
\end{table}

Las variables se agrupan en familias semánticas:

\begin{itemize}
  \item \textbf{Montos y producto}: \texttt{TransactionAmt}, \texttt{ProductCD} (W, C, R, H, S).
  \item \textbf{Tarjeta}: \texttt{card1}--\texttt{card6} (emisor, tipo, categoría).
  \item \textbf{Geografía}: \texttt{addr1}--\texttt{addr2}, \texttt{dist1}--\texttt{dist2}.
  \item \textbf{Correo}: \texttt{P\_emaildomain}, \texttt{R\_emaildomain}.
  \item \textbf{Conteos} \texttt{C1}--\texttt{C14}: cuántos elementos asociados (tarjetas,
        direcciones, correos) aparecen en la ventana de observación.
  \item \textbf{Deltas de tiempo} \texttt{D1}--\texttt{D15}: días transcurridos desde un evento
        previo (p.\,ej.\ primera transacción, cambio de tarjeta).
  \item \textbf{Coincidencias} \texttt{M1}--\texttt{M9}: variables booleanas de matching.
  \item \textbf{\texttt{V1}--\texttt{V339}}: variables propietarias de Vesta (rangos, conteos,
        hashes), con bloques fuertemente correlacionados.
  \item \textbf{Identidad} \texttt{id\_01}--\texttt{id\_38}, \texttt{DeviceType},
        \texttt{DeviceInfo}: huella de dispositivo y de red.
\end{itemize}

\section{Desbalance del problema}

De las 590\,540 transacciones de entrenamiento, 20\,663 son fraudulentas: un 3{,}50\,\%
(razón de desbalance 1:27, \cref{fig:target}). Esta proporción hace que la exactitud
(\emph{accuracy}) sea una métrica inútil ---un clasificador trivial que predice ``legítima''
alcanza 96{,}5\,\% de exactitud con utilidad nula---, por lo que la competencia adoptó el
área bajo la curva ROC (AUC) como métrica oficial, que ordena transacciones en lugar de
etiquetarlas con umbral fijo.

\begin{figure}[ht]
  \centering
  \includegraphics[width=0.62\textwidth]{\figdir fig_distribucion_target.pdf}
  \caption{Distribución de la variable objetivo. El desbalance 1:27 exige métricas de orden
           (AUC) y hace innecesario el re-muestreo agresivo para modelos de árboles.}
  \label{fig:target}
\end{figure}

\section{Calidad de los datos}

Dos problemas dominan la calidad de los datos. Primero, los \textbf{valores faltantes}: la
familia \texttt{V} promedia un 43\,\% de faltantes y la tabla de identidad, al fusionarse por
completitud izquierda, introduce un 76\,\% de faltantes estructurales en sus columnas
(\cref{fig:nulos}). Segundo, la \textbf{cola pesada} del monto (\cref{fig:amt}), que en escala
lineal concentra el 99\,\% de la masa por debajo de \$300 mientras la cola llega a \$31\,937:
los montos extremos son justamente sospechosos, por lo que no deben recortarse sino
re-escalarse (logaritmo) o dejarse en bruto para los árboles.

\begin{figure}[ht]
  \centering
  \includegraphics[width=0.85\textwidth]{\figdir fig_nulos_por_grupo.pdf}
  \caption{Porcentaje promedio de valores faltantes por familia de variables. La ausencia de la
           tabla de identidad es estructural (transacciones sin huella de dispositivo) y, como se
           verá, codificar la ausencia como categoría propia (-999) es preferible a imputar.}
  \label{fig:nulos}
\end{figure}

\begin{figure}[ht]
  \centering
  \includegraphics[width=0.85\textwidth]{\figdir fig_amt_distribucion.pdf}
  \caption{Distribución del monto por clase en escala logarítmica. Las transacciones fraudulentas
           se concentran en montos bajos-repetitivos, con una cola superior más densa que la de
           las legítimas.}
  \label{fig:amt}
\end{figure}

\section{La estructura temporal: la dificultad central}

\texttt{TransactionDT} es un desplazamiento relativo en segundos (693\,999--1\,581\,113\,131)
que cubre 182 días de entrenamiento. La partición oficial de la competencia es
\textbf{temporal}: el conjunto de prueba corresponde al mes siguiente al final del
entrenamiento. Esto tiene tres consecuencias metodológicas:

\begin{enumerate}
  \item Una validación aleatoria (KFold clásico) \emph{sobreestima} el desempeño: filas
        ``del futuro'' aparecen en el entrenamiento de cada fold.
  \item Cualquier variable cuya distribución derive con el tiempo produce una relación espuria
        con el target (\cref{fig:fraude-temporal} y \cref{sec:dcols}).
  \item La validación honesta debe simular el futuro: \emph{holdout} por corte temporal,
        \texttt{TimeSeriesSplit}, o \texttt{GroupKFold} agrupando por mes calendario.
\end{enumerate}

\begin{figure}[ht]
  \centering
  \includegraphics[width=0.95\textwidth]{\figdir fig_fraude_temporal.pdf}
  \caption{Tasa diaria de fraude durante los 182 días de entrenamiento. Los picos agudos
           corresponden a campañas de ataque acotadas en el tiempo: el patrón no es estacionario,
           y por eso la selección de variables por \emph{consistencia temporal}
           (\cref{sec:timeconsistency}) resulta decisiva.}
  \label{fig:fraude-temporal}
\end{figure}

Como referencia de mercado, \cref{fig:leaderboard} muestra la distribución de puntajes del
tablero público (6\,355 equipos): la masa de competidores se concentra bajo 0{,}95 y el top-5
supera 0{,}966, lo que ilustra el diferencial que aportan las técnicas que este informe estudia.

\begin{figure}[ht]
  \centering
  \includegraphics[width=0.9\textwidth]{\figdir fig_leaderboard.pdf}
  \caption{Histograma del tablero público de la competencia. En rojo, los cinco primeros
           equipos: AlKo (0{,}9681), FraudSquad (0{,}9677), Young for you (0{,}9676),
           2 uncles and 3 puppies (0{,}9672) y Mr Lonely (0{,}9663).}
  \label{fig:leaderboard}
\end{figure}
